\begingroup
\setlength{\abovedisplayskip}{7pt}
\setlength{\belowdisplayskip}{7pt}
\setlength{\abovedisplayshortskip}{3pt}
\setlength{\belowdisplayshortskip}{4pt}
\def\ixrad{\mathord{\vcenter{\hbox{\Large\(\surd\)}}}}
\def\ixangle#1#2#3{%
  \makebox[1.95em][r]{\(#1\)}%
  \makebox[1.45em][r]{\(#2\)}%
  \makebox[3.10em][r]{\(#3\)}}
\def\ixobsCount#1#2{\makebox[.8em][r]{#1}\enspace\makebox[2.6em][c]{#2}}
\def\ixobs#1#2#3#4#5#6#7#8#9{%
  #1 & \(\left\{\begin{array}{@{}l@{\enspace}l@{}}
  \text{Direct}&\ixangle{#2}{#3}{#4}\\
  \text{Reflectirt}&\ixangle{#5}{#6}{#7}
  \end{array}\right.\) &
  \begin{tabular}{@{}l@{}}#8\\#9\end{tabular}}

% Printed 32; principal PDF 39; Google PDF 43.
\par\clearpage
\noindent\makebox[12mm][l]{\fontsize{13}{16}\selectfont 32}\parbox[t]{\dimexpr\linewidth-24mm\relax}{\centering\scriptsize\bfseries BESTIMMUNG DES BREITENUNTERSCHIEDES ZWISCHEN GÖTTINGEN UND ALTONA.}\makebox[12mm][r]{}\par\medskip

{\small
\setlength{\arraycolsep}{.85em}
\[
\begin{array}{r||r|r||r|r}
\smash{\raisebox{-.7ex}{Stern}}&
\multicolumn{2}{c||}{\text{\normalsize\textls[100]{Göttingen.}}}&
\multicolumn{2}{c}{\text{\normalsize\textls[100]{Altona.}}}\\
&\text{Coll. F.}&\text{Gewicht}&\text{Coll. F.}&\text{Gewicht}\\
\hline
24&3''{,}93&10{,}91&1''{,}21&12{,}00\\
25&3{,}88&12{,}00&0{,}83&12{,}00\\
26&3{,}54&12{,}00&0{,}76&12{,}00\\
27&4{,}44&4{,}00&1{,}97&7{,}50\\
28&3{,}27&10{,}91&0{,}95&12{,}00\\
29&3{,}89&6{,}00&1{,}15&12{,}00\\
30&4{,}54&10{,}91&1{,}70&12{,}00\\
31&4{,}52&10{,}91&1{,}73&12{,}00\\
32&3{,}26&12{,}00&1{,}13&12{,}00\\
33&3{,}49&10{,}91&1{,}11&10{,}91\\
34&4{,}58&12{,}00&1{,}66&12{,}00\\
35&3{,}61&10{,}00&1{,}37&10{,}00\\
36&4{,}33&12{,}00&1{,}34&10{,}00\\
37&3{,}97&12{,}00&1{,}91&10{,}91\\
38&3{,}52&12{,}00&1{,}56&10{,}00\\
39&4{,}16&3{,}00&1{,}45&8{,}00\\
40&5{,}04&4{,}80&1{,}02&6{,}00\\
41&3{,}46&4{,}80&1{,}75&10{,}00\\
42&3{,}84&4{,}00&2{,}45&8{,}89\\
43&4{,}40&4{,}00&1{,}01&10{,}00
\end{array}
\]
\par
}

Die Mittelwerthe sind folgende:
\begin{center}
\begin{tabular}{@{}l@{\qquad}r@{\enspace}l@{}}
Collimationsfehler in Göttingen & \(3''{,}75\) & mit dem Gewicht \(455{,}17\)\\
Collimationsfehler in Altona & \(1{,}40\) & mit dem Gewicht \(432{,}18\).
\end{tabular}
\end{center}

Die Realität der Veränderung des Collimationsfehlers ist offenbar, und
es leidet keinen Zweifel, dass dieselbe auf dem obwohl mit aller möglichen
Vorsicht geleiteten Transport eingetreten ist.

\begin{center}\textbf{4.}\end{center}

Obgleich man sich bei dem für den Breitenunterschied gefundenen
Resultate vollkommen beruhigen kann, so ist es doch wenigstens in theoretischer
Rücksicht nicht überflüssig zu bemerken, dass die im 1. Art. angewandte
Combination der Beobachtungen noch nicht die möglich vortheilhafteste ist,
insofern nicht an jedem Ort jeder Stern in der einen Lage des Sectors eben so
oft beobachtet ist, wie in der andern. In der That hat die Bestimmung der
wahren Zenithdistanz in Göttingen durch die Formel \(\tfrac12(a+a')\) das Gewicht
\(\dfrac{4\alpha\alpha'}{\alpha+\alpha'}\); wäre nun der Collimationsfehler in Göttingen genau
bekannt und \(=f\),
so würde die Bestimmung der wahren Zenithdistanz daselbst durch die Formel
\[
\frac{\alpha(a+f)+\alpha'(a'-f)}{\alpha+\alpha'}
\]

% Printed 33; principal PDF 40; Google PDF 44.
\par\clearpage
\noindent\makebox[12mm][l]{}\makebox[\dimexpr\linewidth-24mm\relax][c]{\small\textbf{III.\quad RESULTATE.}}\makebox[12mm][r]{\fontsize{13}{16}\selectfont 33}\par\medskip

\noindent das Gewicht \(\alpha+\alpha'=\dfrac{4\alpha\alpha'}{\alpha+\alpha'}+
\dfrac{(\alpha-\alpha')^2}{\alpha+\alpha'}\) haben, d.\ i. ein grösseres als nach der
andern Methode, so oft \(\alpha\) und \(\alpha'\) ungleich sind. Eben so verhält es sich mit
der wahren Zenithdistanz in Altona, und auf diese Art würden selbst einseitige
Beobachtungen (wie die von Nro. 5) einen, wenn auch nur geringen,
Beitrag zur Vergrösserung der Genauigkeit geben. Nun sind zwar die Collimationsfehler
an beiden Plätzen nicht mit \textls[100]{absoluter} Schärfe bekannt: allein
man überzeugt sich leicht, dass die Anwendung der für dieselben gefundenen
Mittelwerthe das Gewicht nur ganz unbedeutend vermindert.

\begin{center}\textbf{5.}\end{center}

Will man jedoch ein reines, den Forderungen der strengen Theorie ganz
Genüge leistendes Resultat erhalten, so muss man die Bestimmung des
Breitenunterschiedes, der Collimationsfehler und der wahren Zenithdistanzen der
einzelnen Sterne an dem einen Ort als Ein Problem behandeln, wo diese
unbekannten Grössen (in unserm Fall 46 an der Zahl) aus den sämmtlichen
durch sie bestimmten beobachteten Grössen (171) durch eben so viele
Gleichungen abgeleitet werden müssen, indem diese nach den Vorschriften der
Wahrscheinlichkeitsrechnung combinirt werden. Setzt man die Collimationsfehler
in Göttingen und Altona \(=f\) und \(g\), den Breitenunterschied \(=h\), die
wahre Zenithdistanz eines Sterns in Göttingen \(=k\), so hat man aus den
Beobachtungen dieses Sterns die vier Gleichungen mit den Gewichten \(\alpha,\alpha',\beta,\beta'\):
\[
\begin{aligned}
a&=k-f\\
a'&=k+f\\
b&=k-g-h\\
b'&=k+g-h.
\end{aligned}
\]

Es ist kaum nöthig zu erinnern, dass es zur Erleichterung der Rechnung
vortheilhafter ist, anstatt jener unbekannten Grössen, die noch erforderlichen
Correctionen einzuführen, welche an die schon sehr nahe bestimmten Werthe
anzubringen sind; lassen wir die Zeichen \(f^0,g^0,h^0,k^0\) diese genäherten Werthe
bedeuten, so mag man annehmen
\[
k^0=\frac{\alpha(a+f^0)+\alpha'(a'-f^0)+\beta(b+g^0+h^0)+\beta'(b'-g^0+h^0)}
{\alpha+\alpha'+\beta+\beta'}.
\]

Bei Befolgung jener Vorschrift (welche man bei Anwendung der Methode
\par\smallskip
\noindent\textbf{IX.}\hfill 5\hspace{.12\textwidth}\par

% Printed 34; principal PDF 41; Google PDF 45.
\par\clearpage
\noindent\makebox[12mm][l]{\fontsize{13}{16}\selectfont 34}\parbox[t]{\dimexpr\linewidth-24mm\relax}{\centering\scriptsize\bfseries BESTIMMUNG DES BREITENUNTERSCHIEDES ZWISCHEN GÖTTINGEN UND ALTONA.}\makebox[12mm][r]{}\par\medskip

\noindent der kleinsten Quadrate auf nur etwas zusammengesetzte Fälle niemals aus
den Augen setzen sollte) und dem Gebrauch einer schicklichen \textls[100]{indirecten}
Auflösungsmethode verwandelt sich eine Arbeit, die ohne jene und bei
directer Elimination unerträglich weitläuftig ausfällt, in ein leichtes Spiel.

\begin{center}\textbf{6.}\end{center}

Der Erfolg dieser Rechnung, welche ausführlich herzusetzen unnöthig
wäre, ist, dass die frühern Bestimmungen gar keine merkliche Correction
erhalten. Es findet sich die Verbesserung des Breitenunterschiedes \(=-0''{,}014\),
die Verbesserung des Collimationsfehlers in Göttingen \(=+0''{,}012\), die
Verbesserung des Collimationsfehlers in Altona \(=-0''{,}014\); folglich die neuen
Bestimmungen
\begin{center}
\begin{tabular}{@{}l@{\quad}r@{\enspace}c@{\enspace}l@{}}
Breitenunterschied & \(2^\circ0'56''{,}51\), & Gewicht & \(=217{,}67\)\\
Collimationsfehler in Göttingen & \(3{,}76\) & » & \(\phantom{=}457{,}03\)\\
Collimationsfehler in Altona & \(1{,}39\) & » & \(\phantom{=}437{,}64\).
\end{tabular}
\end{center}

Die Veränderungen der nach der Vorschrift des vorhergehenden Artikels
zum Grunde gelegten wahren Zenithdistanzen der einzelnen Sterne in Göttingen
sind gleichfalls fast alle unter \(0''{,}01\). Die sich ergebenden Werthe hier
aufzuführen, wäre überflüssig, da es dieselben sind, aus welchen die oben
mitgetheilten Declinationen der Sterne unter Voraussetzung der Polhöhe des
Beobachtungsplatzes \(51^\circ31'47''{,}92\) abgeleitet sind. Dagegen setzen wir die
Unterschiede hier her, welche nach Substitution der gefundenen Werthe in den
171 Gleichungen übrig bleiben.

{\small
\[
\begin{array}{@{}c|c||c|c||c|c||c|c@{}}
\text{Stern}&\text{Unterschied}&\text{Stern}&\text{Unterschied}&
\text{Stern}&\text{Unterschied}&\text{Stern}&\text{Unterschied}\\[.8ex]
1&+0''{,}07&5&+0''{,}03&9&-0''{,}06&13&+0''{,}87\\
&+0{,}09&&-0{,}03&&-0{,}09&&+0{,}30\\
&-0{,}26&&\text{—}&&-0{,}23&&-1{,}35\\
&+0{,}13&&-0{,}03&&+0{,}97&&-0{,}04\\
2&+0{,}56&6&+0{,}62&10&-0{,}71&14&+0{,}40\\
&-0{,}08&&+0{,}03&&-0{,}25&&+0{,}21\\
&-0{,}12&&-1{,}95&&+0{,}70&&-0{,}29\\
&-0{,}53&&-0{,}02&&+0{,}44&&-0{,}48\\
3&+0{,}48&7&-0{,}52&11&+0{,}12&15&-0{,}04\\
&+0{,}34&&+0{,}16&&-1{,}03&&+0{,}19\\
&-0{,}70&&-0{,}08&&+0{,}72&&-0{,}04\\
&-0{,}18&&+0{,}52&&+0{,}71&&-0{,}11\\
4&-0{,}35&8&-0{,}56&12&+0{,}52&16&-0{,}07\\
&-0{,}35&&+0{,}12&&-0{,}87&&-0{,}60\\
&+0{,}79&&+0{,}76&&+0{,}80&&-0{,}05\\
&-0{,}17&&-0{,}23&&-0{,}24&&+0{,}72
\end{array}
\]
\par
}

% Printed 35; principal PDF 42; Google PDF 46.
\par\clearpage
\noindent\makebox[12mm][l]{}\makebox[\dimexpr\linewidth-24mm\relax][c]{\small\textbf{III.\quad RESULTATE.}}\makebox[12mm][r]{\fontsize{13}{16}\selectfont 35}\par\medskip

{\small
\[
\begin{array}{@{}c|c||c|c||c|c||c|c@{}}
\text{Stern}&\text{Unterschied}&\text{Stern}&\text{Unterschied}&
\text{Stern}&\text{Unterschied}&\text{Stern}&\text{Unterschied}\\[.8ex]
17&-0''{,}06&24&-0''{,}17&31&-0''{,}41&38&+0''{,}30\\
&-0{,}52&&+0{,}18&&+1{,}11&&-0{,}19\\
&+0{,}96&&+0{,}16&&-0{,}59&&-0{,}24\\
&-0{,}35&&-0{,}20&&+0{,}09&&+0{,}11\\
18&-0{,}23&25&-0{,}03&32&+0{,}14&39&+0{,}90\\
&-0{,}49&&+0{,}22&&-0{,}87&&+1{,}71\\
&+1{,}09&&+0{,}46&&+0{,}63&&-0{,}82\\
&-0{,}34&&-0{,}67&&+0{,}11&&-0{,}69\\
19&+0{,}32&26&+0{,}90&33&+0{,}19&40&-1{,}81\\
&-0{,}36&&+0{,}46&&-0{,}35&&+0{,}75\\
&+0{,}10&&-0{,}06&&+0{,}41&&+0{,}60\\
&+0{,}08&&-1{,}32&&-0{,}16&&-0{,}14\\
20&-0{,}06&27&-0{,}14&34&-0{,}48&41&+0{,}39\\
&-0{,}26&&+1{,}22&&+1{,}17&&-0{,}21\\
&+0{,}66&&-0{,}71&&-0{,}62&&-0{,}38\\
&-0{,}44&&+0{,}46&&-0{,}07&&+0{,}35\\
21&-0{,}22&28&+0{,}46&35&-0{,}08&42&+0{,}09\\
&-0{,}85&&-0{,}53&&-0{,}37&&+0{,}25\\
&+0{,}63&&+0{,}52&&+0{,}25&&-1{,}02\\
&+0{,}37&&-0{,}36&&+0{,}21&&+1{,}11\\
22&+0{,}77&29&-0{,}80&36&-0{,}14&43&-0{,}42\\
&-0{,}55&&-0{,}53&&+1{,}02&&+0{,}86\\
&-1{,}55&&+0{,}58&&-0{,}47&&+0{,}29\\
&+1{,}60&&+0{,}10&&-0{,}58&&-0{,}47\\
23&-0{,}80&30&-0{,}83&37&+0{,}11&&\\
&+0{,}67&&+0{,}74&&+0{,}54&&\\
&-0{,}03&&-0{,}21&&-0{,}93&&\\
&+0{,}03&&+0{,}42&&+0{,}11&&
\end{array}
\]
\par
}

\begin{center}\textbf{7.}\end{center}

Die Summe der Producte aus den Quadraten dieser 171 Unterschiede in
die entsprechende Anzahl der Beobachtungen findet sich \(=292{,}8249\). Nach
dem bereits angeführten Lehrsatz (Theoria Comb. Observ. Art. 38) hat man
als genäherten Werth des mittlern Fehlers einer einfachen Beobachtung die
Quadratwurzel aus dem Bruch zu betrachten, dessen Zähler jene Summe, und
dessen Nenner der Überschuss der Anzahl der verglichenen Beobachtungsdata
über die Anzahl der nach der Methode der kleinsten Quadrate daraus
abgeleiteten unbekannten Grössen ist, in unserm Falle \(171-46=125\). Es findet
sich hieraus jener mittlere Fehler \(=1''{,}5308\), wenig von dem im 2. Art.
gefundenen verschieden. Der mittlere in dem Endresultate für den
Breitenunterschied zu befürchtende Fehler würde demnach sein
\[
=\frac{1''{,}5308}{\surd 217{,}67}=0''{,}1038.
\]
\par\vfill
\noindent\hfill \(5^{*}\)\hspace{.12\textwidth}\par

% Printed 36; principal PDF 43; Google PDF 47.
\par\clearpage
\noindent\makebox[12mm][l]{\fontsize{13}{16}\selectfont 36}\parbox[t]{\dimexpr\linewidth-24mm\relax}{\centering\scriptsize\bfseries BESTIMMUNG DES BREITENUNTERSCHIEDES ZWISCHEN GÖTTINGEN UND ALTONA.}\makebox[12mm][r]{}\par\medskip

\begin{center}\textbf{8.}\end{center}

Bei den bisherigen Rechnungen ist vorausgesetzt, dass alle den verschiedenen
Beobachtungen anhängenden Fehler als völlig unabhängig von einander
oder als rein zufällig betrachtet werden können. Diese Voraussetzung aber
ist offenbar nicht ganz richtig, indem alle \(\alpha\) Beobachtungen, welche zu der
Bestimmung eines \(a\) concurrirt haben, nach der Natur des Instruments sich
auf einen und denselben Theilungspunkt beziehen, und also ausser den
eigentlichen rein zufälligen Beobachtungsfehlern noch den Fehler der Theilung bei
diesem Punkte involviren. Dasselbe gilt von \(a'\), \(b\) und \(b'\). Die
Theilungsfehler sind ihrerseits unbekannte Grössen, die in Beziehung auf die einzelnen
171 Beobachtungsresultate auch als rein zufällig und von einander unabhängig
betrachtet werden mögen, da man die Fälle, wo verschiedene derselben sich
auf einerlei Theilungspunkt bezogen haben, ihrer geringen Anzahl wegen
ignoriren kann. Die Berücksichtigung dieses Umstandes macht nun eine
Modification obiger Rechnungen nothwendig, obwohl am Ende in praktischer
Rücksicht die Resultate gar nicht geändert werden.

Bezeichnet man den eigentlichen mittlern Beobachtungsfehler, der nur
von zufälligen Ursachen mit Ausschluss der Theilungsfehler herrührt, mit \(m\),
und den mittlern Theilungsfehler mit \(\mu\), so wird der vollständige mittlere
Beobachtungsfehler \(=\surd(mm+\mu\mu)\) zu setzen sein, und der mittlere Fehler eines
Mittels aus \(\alpha\) Beobachtungen, die sich auf einerlei Theilungspunkt beziehen,
\[
=\ixrad\left(\frac{mm}{\alpha}+\mu\mu\right),
\]
\noindent oder wenn wir \(\mu\mu=mm\theta\) setzen,
\[
=m\ixrad\left(\frac1\alpha+\theta\right).
\]

Insofern wir also das Gewicht einer Beobachtung, \textls[100]{ohne Theilungsfehler},
zur Einheit annehmen, wird das Gewicht von \(a\) nunmehro
\[
=\frac{\alpha}{1+\alpha\theta}
\]
\noindent sein, und eben so die Gewichte von \(a'\), \(b\), \(b'\) resp.
\[
=\frac{\alpha'}{1+\alpha'\theta},\qquad
\frac{\beta}{1+\beta\theta},\qquad
\frac{\beta'}{1+\beta'\theta}.
\]
Bei der erstern Combinationsmethode wird man daher das Gewicht des
\par

% Printed 37; principal PDF 44; Google PDF 48.
\par\clearpage
\noindent\makebox[12mm][l]{}\makebox[\dimexpr\linewidth-24mm\relax][c]{\small\textbf{III.\quad RESULTATE.}}\makebox[12mm][r]{\fontsize{13}{16}\selectfont 37}\par\medskip

\noindent Resultats für den Breitenunterschied aus den Beobachtungen eines Sterns,
wenn man den vorigen Ausdruck
\[
\frac4{\dfrac1\alpha+\dfrac1{\alpha'}+\dfrac1\beta+\dfrac1{\beta'}}=p
\]
\noindent setzt, jetzt
\[
=\frac p{1+p\theta}
\]
\noindent zu setzen, und nach Maassgabe dieser Gewichte aus den 42 Bestimmungen
das Mittel zu nehmen haben. Bei der zweiten Combinationsmethode hingegen
hat man nur jeder der 171 Gleichungen ein Gewicht beizulegen, welches
durch eine der Formeln \(\dfrac{\alpha}{1+\alpha\theta}\) u.\ s.\ w. bestimmt wird.

Offenbar kann eine Veränderung des Endresultats selbst sowohl, als des
mittlern in demselben zu befürchtenden Fehlers nur dadurch eintreten, dass
die neuen Gewichte den frühern nicht proportional sind. Bei der vorigen
Methode waren nur die Resultate der zahlreichern Beobachtungsreihen etwas
zu viel bevorzugt; die Berücksichtigung der Theilungsfehler bringt ihre
Gewichte der Gleichheit näher, desto mehr, je grösser die Theilungsfehler
vorausgesetzt werden, so dass bei Beobachtungen mit einem Instrumente, wo die
Theilungsfehler die eigentlichen Beobachtungsfehler sehr weit überwögen, man
sich nur begnügen könnte, alle Bestimmungen als gleich zuverlässig zu
betrachten.

\begin{center}\textbf{9.}\end{center}

Die angezeigten Methoden haben also gar keine Schwierigkeit, sobald nur
der Coefficient \(\theta\) bekannt ist. Man kann zu einer genäherten Kenntniss desselben,
auf welche es hier begreiflich nur ankommt, auf einem indirecten
Wege gelangen.

Wir bemerken zuvörderst, dass die Beobachtungen selbst ein Mittel darbieten,
den eigentlichen mittlern Beobachtungsfehler \(m\) mit sehr grosser
Zuverlässigkeit zu bestimmen. In der That macht sich derselbe unabhängig von
dem Theilungsfehler in den Unterschieden der \textls[100]{einzelnen} Werthe, aus denen
jedes \(a\) (oder \(a'\), \(b\), \(b'\)) das Mittel ist, von einander oder von diesem Mittel,
bemerkbar, und wenn \(\alpha\) sehr gross wäre, so würde die Summe der Quadrate
dieser Unterschiede der einzelnen Werthe von \(a\) vom Mittel als eine ge-
\par

% Printed 38; principal PDF 45; Google PDF 49.
\par\clearpage
\noindent\makebox[12mm][l]{\fontsize{13}{16}\selectfont 38}\parbox[t]{\dimexpr\linewidth-24mm\relax}{\centering\scriptsize\bfseries BESTIMMUNG DES BREITENUNTERSCHIEDES ZWISCHEN GÖTTINGEN UND ALTONA.}\makebox[12mm][r]{}\par\medskip

\noindent näherte Bestimmung von \((\alpha-1)mm\) anzusehen sein. Eine solche einzelne
Bestimmung kann nun zwar in unserm Fall, wo \(\alpha\) nie grösser als 7 ist, von
dem richtigen Werthe sehr abweichen; allein die Summe aller 171 partiellen
Summen (für alle \(a\), \(a'\), \(b\), \(b'\) und für alle Sterne) muss nach den Grundsätzen
der Wahrscheinlichkeitsrechnung von
\[
(\Sigma(\alpha-1)+\Sigma(\alpha'-1)+\Sigma(\beta-1)+\Sigma(\beta'-1))mm,
\]
\noindent in unserm Fall von \(728mm\), wenig verschieden sein. Wir haben jene Summe
der 171 partiellen Summen
\[
=844{,}50
\]
\noindent gefunden, woraus sich für \(m\) der sehr zuverlässige Werth
\[
=1''{,}0770
\]
\noindent ergibt, bedeutend kleiner als der im 2. und 7. Artikel gefundene. Es bestätigt
sich also die Einwirkung der Theilungsfehler vollkommen, um derenwillen
die früher herausgebrachten Zahlen kein reines Resultat geben konnten.

\begin{center}\textbf{10.}\end{center}

In Ermangelung einer directen Kenntniss des mittlern Theilungsfehlers
kann man nun \(\theta\) auf eine indirecte Art so bestimmen, dass beim Gebrauch
der ersten Methode, nach dem Verfahren des Art. 2, oder beim Gebrauch
der zweiten Methode nach dem Verfahren des Art. 7, der mittlere Fehler
einer Beobachtung, deren Gewicht als Einheit angenommen war, wiederum
dem gefundenen Werthe von \(m\) gleich wird.

Es hat indessen nicht belohnend genug geschienen, solche Versuche so
lange zu wiederholen, bis eine vollkommene Übereinstimmung erreicht wäre.
Vielmehr schien es hinreichend, nachdem durch anderweitige Betrachtungen
erkannt war, dass der letzte Werth von \(\theta\) nur wenig von \(0{,}2\) verschieden
ausfallen könnte, diesen Werth bloss der ersten Combinationsmethode unterzulegen,
woraus sich dann ergeben hat
\begin{center}
\begin{tabular}{@{}l@{\quad}l@{}}
Breitenunterschied & \(=2^\circ0'56''{,}50\)\\
Gewicht dieser Bestimmung & \(=104{,}29\)\\
\multicolumn{2}{@{}l@{}}{Mittlerer Fehler einer Beobachtung, deren Gewicht die Einheit,}
\end{tabular}
\end{center}
\[
=1''{,}131
\]

% Printed 39; principal PDF 46; Google PDF 50.
\par\clearpage
\noindent\makebox[12mm][l]{}\makebox[\dimexpr\linewidth-24mm\relax][c]{\small\textbf{III.\quad RESULTATE.}}\makebox[12mm][r]{\fontsize{13}{16}\selectfont 39}\par\medskip

\noindent und daher der mittlere in obigem Endresultate zu befürchtende Fehler
\[
=0''{,}1108.
\]

Die Anwendung der zweiten Combinationsmethode, mit demselben Werthe
von \(\theta\), würde vermuthlich eine noch nähere Übereinstimmung mit obigem
Werthe von \(m\) hervorgebracht, das Endresultat für den Breitenunterschied
vielleicht um \(0''{,}01\) vermindert, das Gewicht dieser Bestimmung gewiss etwas
weniges vergrössert haben; es wurde aber der Mühe nicht werth gehalten,
deshalb diese Rechnung von neuem durchzuführen. Man kann sich also an
den gefundenen Breitenunterschied \(2^\circ5'56''50\) halten, und dessen Fehler als
wahrscheinlich zwischen den Grenzen \(\pm0''{,}07\) enthalten ansehen.

\begin{center}\textbf{11.}\end{center}

Wenn wir den obigen Werth von \(\theta\) beibehalten, so ergibt sich der mittlere
Theilungsfehler \(=m\surd\theta=0''{,}48\), daher der sogenannte wahrscheinliche
Theilungsfehler der einzelnen Punkte \(=0''{,}32\) gesetzt werden mag. Offenbar
bezieht sich dies aber nur auf die unregelmässigen Theilungsfehler, oder auf
die Abweichungen der einzelnen Punkte von einer fingirten, sich diesen so
genau wie möglich anschliessenden gleichförmigen Theilung, deren
\textls[100]{absolute} Richtigkeit hiebei eigentlich gar nicht in Frage kommen konnte.
Oder mit andern Worten, das gefundene Resultat für den Breitenunterschied mit
der ihm beigelegten Genauigkeit bezieht sich, streng genommen, nur auf mittlere
Sectorgrade, und bleibt von der absoluten Richtigkeit derselben abhängig.
Dem Astronomen bietet das Instrument gar kein selbstständiges Mittel dar,
\textls[100]{diese} zu prüfen. Wenn man indessen erwägt, dass die Endpunkte des Bogens
von dem Künstler mit äusserster Sorgfalt niedergelegt sind, und dass hier nur
von einem kleinen Theile des ganzen Bogens die Rede ist, so wird man
zugeben müssen, dass die Unsicherheit des gefundenen Breitenunterschiedes aus
dieser Quelle nur um ein sehr Geringes vergrössert werden kann. Einige
Controlle für die absolute Richtigkeit der Theilung geben übrigens auch die
von mir am \textsc{Reichenbach}schen Meridiankreise beobachteten Zenithdistanzen
derselben 43 Sterne, deren Unterschiede von den am Sector beobachteten, bei
einer Anordnung nach den Declinationen, keine Spur von Regelmässigkeit
zeigen.

% Printed 40; principal PDF 47; Google PDF 51.
\par\clearpage
\noindent\makebox[12mm][l]{\fontsize{13}{16}\selectfont 40}\parbox[t]{\dimexpr\linewidth-24mm\relax}{\centering\scriptsize\bfseries BESTIMMUNG DES BREITENUNTERSCHIEDES ZWISCHEN GÖTTINGEN UND ALTONA.}\makebox[12mm][r]{}\par\medskip

\begin{center}\textbf{12.}\end{center}

Der Platz des Mittelpunkts des Sectors in Göttingen war \(1{,}060\) Toisen
nördlich und \(7{,}595\) Toisen östlich vom Centrum der Axe des \textsc{Reichenbach}schen
Meridiankreises; in Altona hingegen war der Mittelpunkt des Sectors \(13{,}511\)
Toisen südlich und \(2{,}578\) westlich vom Mittelpunkt des dortigen
Meridiankreises. Die Reduction des Breitenunterschiedes der Sectorplätze auf den der
Meridiankreise ist daher für Göttingen \(0''{,}07\) und für Altona \(0''{,}85\), und folglich
der Breitenunterschied der Sternwarten von Göttingen und Altona in
Beziehung auf die Plätze der \textsc{Reichenbach}schen Meridiankreise
\[
=2^\circ0'57''{,}42.
\]

\begin{center}\textbf{13.}\end{center}

Die absolute Polhöhe, welche den oben gegebenen aus den Zenithdistanzen
abgeleiteten Declinationen der Sterne zum Grunde gelegt ist, beruht auf 89
Beobachtungen des Nordsterns, am \textsc{Reichenbach}schen Meridiankreise, in beiden
Culminationen, direct und von einer Wasserfläche reflectirt. Da die
Beobachtungen von 1824, welche den grössten Theil ausmachen, bisher noch nicht
bekannt gemacht sind, so stelle ich hier sämmtliche Beobachtungen zusammen,
und bemerke nur, dass meistens die directe Einstellung beim Antritt an den
zweiten, vierten (mittelsten) und sechsten Faden, die Einstellung des
reflectirten Bildes hingegen beim Antritt an den ersten, dritten, fünften und
siebenten Faden gemacht ist. Von diesen auf die Culminationszeit reducirten
Zenithdistanzen ist hier das Mittel angegeben, welches bloss von der
Refraction nach \textsc{Bessel}s Tafeln befreit ist, also Collimationsfehler und Wirkung
der Biegung des Fernrohrs noch einschliesst.

\begin{center}
\textls[100]{Zenithdistanzen des Nordsterns.}\par\smallskip
1820.\quad Kreis in Osten.\par\smallskip
\renewcommand{\arraystretch}{1.12}%
\begin{tabular}{@{}r@{\enspace}l|l@{}}
\ixobs{Mai 13. Untere Culm.}{319^\circ}{50'}{20''{,}73}{220}{5}{3{,}94}{\ixobsCount{3}{Beob.}}{\ixobsCount{4}{»}}\\[.6ex]
\ixobs{13. Obere Culm.}{323}{8}{41{,}51}{216}{46}{44{,}31}{\ixobsCount{1}{»}}{\ixobsCount{1}{»}}
\end{tabular}
\end{center}

% Printed 41; principal PDF 48; Google PDF 52.
\par\clearpage
\noindent\makebox[12mm][l]{}\makebox[\dimexpr\linewidth-24mm\relax][c]{\small\textbf{III.\quad RESULTATE.}}\makebox[12mm][r]{\fontsize{13}{16}\selectfont 41}\par\medskip

\begin{center}
1824.\quad Kreis in Osten.\par\smallskip
\renewcommand{\arraystretch}{1.18}%
\begin{tabular}{@{}r@{\enspace}l|l@{}}
\ixobs{Apr, 20. Obere Culm.}{323^\circ}{7'}{52''{,}62}{216}{48}{54{,}93}{\ixobsCount{1}{Beob.}}{\ixobsCount{2}{»}}\\[.7ex]
\ixobs{21. Untere Culm.}{319}{52}{30{,}27}{220}{4}{19{,}32}{\ixobsCount{3}{»}}{\ixobsCount{4}{»}}\\[.7ex]
\ixobs{21. Obere Culm.}{323}{7}{54{,}16}{216}{48}{54{,}21}{\ixobsCount{3}{»}}{\ixobsCount{4}{»}}\\[.7ex]
\ixobs{25. Untere Culm.}{319}{52}{30{,}03}{220}{4}{21{,}10}{\ixobsCount{3}{»}}{\ixobsCount{4}{»}}\\[.7ex]
\ixobs{27. Obere Culm.}{323}{7}{55{,}70}{216}{48}{52{,}93}{\ixobsCount{3}{»}}{\ixobsCount{4}{»}}\\[.7ex]
\ixobs{28. Obere Culm.}{323}{7}{55{,}40}{216}{48}{52{,}22}{\ixobsCount{3}{»}}{\ixobsCount{4}{»}}\\[.7ex]
\ixobs{29. Untere Culm.}{319}{52}{29{,}17}{220}{4}{21{,}34}{\ixobsCount{3}{»}}{\ixobsCount{4}{»}}\\[.7ex]
\ixobs{Mai 1. Untere Culm.}{319}{52}{28{,}59}{220}{4}{22{,}62}{\ixobsCount{3}{»}}{\ixobsCount{4}{»}}\\[.7ex]
\ixobs{1. Obere Culm.}{323}{7}{57{,}22}{216}{48}{51{,}66}{\ixobsCount{3}{»}}{\ixobsCount{4}{»}}\\
\multicolumn{2}{@{}c|}{\rule{0pt}{3.3ex}Kreis in Westen.}&\\[.5ex]
\ixobs{Mai 2. Untere Culm.}{40^\circ}{4'}{20''{,}00}{139}{52}{27{,}15}{\ixobsCount{3}{Beob.}}{\ixobsCount{4}{»}}\\[.7ex]
\ixobs{8. Obere Culm.}{36}{48}{49{,}32}{143}{7}{57{,}63}{\ixobsCount{3}{»}}{\ixobsCount{4}{»}}\\[.7ex]
\ixobs{9. Untere Culm.}{40}{4}{22{,}93}{139}{52}{25{,}68}{\ixobsCount{3}{»}}{\ixobsCount{4}{»}}
\end{tabular}
\end{center}

Die Änderungen der Declination des Nordsterns ergeben sich aus \textsc{Bessel}s
Tafeln wie folgt:
\par\smallskip
\noindent\textbf{IX.}\hfill 6\hspace{.12\textwidth}\par
\par
\endgroup
